\documentclass[10pt,a4paper]{article}
\usepackage[margin=25mm]{geometry}
\usepackage{amsmath,amssymb,lmodern}
\setlength{\emergencystretch}{2em}
\begin{document}
\begin{center}
{\Large Finite edge graphs cannot realize every real value}\\[4pt]
{\large Disproof of Conjecture 00000008883}\\[6pt]
ziangni-sys \qquad 4 October 2026
\end{center}
\section*{The finite-realization clause}
Both originals assert that every real surreal can be realized by a
finite edge graph in blue-red Hackenbush. We disprove this clause by a
cardinality obstruction that applies to every possible deterministic
value rule on finite grounded colored graphs. It does not require a
formula for Hackenbush values or an analysis of infinite positions.

A blue-red Hackenbush position is an abstract colored graph with a
ground set: a player removes an edge of the appropriate color, after
which edges disconnected from ground disappear. Only edge incidence,
colors and grounding affect the game. Relabeling vertices or edges and
redrawing the same graph do not change the legal moves or the game
value. Distinct real numbers represent distinct real surreals, so
realizing all real surreals would require realizing all real numbers as
values of these positions.

\section*{Exhaustive finite descriptions}
Allow a grounded multigraph with any finite edge set $E$, any vertex
set $V$, endpoint maps $s,t:E\to V$, edge colors $b:E\to\{0,1\}$,
and ground flags $a:V\to\{0,1\}$. This includes parallel edges and
loops, so we do not obtain a bound by restricting to simple graphs.
The vertex set itself need not be finite. Its incident support is
\[
 S=s(E)\cup t(E),
\]
a finite set. Restricting to $S$ retains every edge and both endpoints,
all edge colors, and the ground status of every incident vertex.
Vertices outside $S$ have no incident edge and can never provide a
move or a connection between edges. Removing them changes no play.
This also handles the source's wording ``finite edges'' without silently
assuming a finite vertex set.

Enumerate $S$ and $E$ as $\{0,\ldots,v-1\}$ and
$\{0,\ldots,e-1\}$. The normalized graph becomes a finite description
consisting of
\[
 s,t:\mathrm{Fin}(e)\to\mathrm{Fin}(v),\quad
 b:\mathrm{Fin}(e)\to\{0,1\},\quad
 a:\mathrm{Fin}(v)\to\{0,1\}.
\]
For fixed $v,e$ there are finitely many such descriptions
($v^{2e}2^{e+v}$, with the usual empty-function conventions).
Taking the union over all natural $v,e$ gives a countable type $D$.
The labeling is exhaustive: any finite-edge position first normalizes
as above and then has one of these descriptions, with all incidence,
color and ground data preserved. Any disconnected or initially
unplayable description only enlarges this countable set.

\section*{No value rule can cover the reals}
The real numbers are uncountable. Hence no function $D\to\mathbb R$
is surjective. Even if some descriptions do not have a real value,
the subtype of descriptions that do is still countable, so no partial
real-value rule can be surjective either.
Equivalently, for any relation $R(d,r)$ in which a description has at
most one real value, the assertion
$\forall r\in\mathbb R\ \exists d\in D\ R(d,r)$ is false.
A Hackenbush position has one game value, and isomorphic descriptions
have the same value. Thus the conjectured realization of every real
surreal by finite edge graphs would give exactly this impossible
coverage. The finite-realization clause is false, irrespective of the
separate claimed correspondence for infinite Hackenbush.

\newpage
\section*{Formal model and completeness}
\raggedright
Lean defines actual graph data \texttt{Graph V E} by endpoint maps,
Boolean edge colors and Boolean ground flags. The edge type is
independent of endpoints, so multiple edges are genuinely allowed.
\texttt{Description v e} contains these same maps over actual finite
index types. The full description type is a dependent sum over all
natural vertex and edge counts, not a bounded collection.
Mathlib's countability handler encodes each fixed-size description,
and countability instances cover the dependent sum.

\texttt{IsRelabelling} explicitly requires both endpoint equations,
the color equation for every edge, and the ground equation for every
vertex, under actual vertex and edge bijections.
\texttt{describe\_preserves\_everything} constructs finite index
representations for arbitrary finite vertex and edge types and proves
all these equations. No completeness of labeling is assumed.

To cover arbitrary vertex sets, \texttt{support} is the actual finite
union of endpoint images. The source/target support theorems prove that
every endpoint lies there. \texttt{normalize} uses the actual subtype
of incident vertices and retains the same edge type and colors.
\texttt{normalization\_preserves\_incidence} proves all endpoint,
color and ground-flag equalities after forgetting the subtype labels.
\texttt{exhaustive\_finite\_edge\_representation} combines this
normalization with finite enumeration. It quantifies over every vertex
type and every finite edge type and produces a finite description plus
actual vertex/edge equivalences satisfying the full relabeling condition.

\section*{Formal cardinal contradiction}
The type \texttt{FiniteDescription} is proved countable.
\texttt{no\_value\_assignment\_realizes\_all\_reals} proves that
\emph{every} function from that type to the actual real numbers is
nonsurjective. \texttt{omitted\_real} gives an omitted real for each
such function. \texttt{no\_partial\_value\_assignment} covers any
subtype of eligible descriptions, so real-valuedness of every position
is not a hidden premise.

The theorem \texttt{no\_functional\_realization\_covers\_all\_reals}
quantifies over any realization relation with at most one real value
per description and negates coverage of all reals. Its proof selects
the unique value on eligible descriptions; full coverage would make
that selected map surjective, contradicting the preceding theorem.
This is a universal necessary-condition disproof: no special
Hackenbush evaluation table, dyadic-value theorem, or desired
nonsurjectivity is assumed. The graph semantics supplies only the
ordinary uniqueness of a position's value and invariance under
relabeling/removal of vertices incident to no edge, justified above.

\section*{Verification}
The project uses Lean 4.19.0 with pinned Mathlib commit
\texttt{c44e0c8ee63ca166450922a373c7409c5d26b00b}.
Axiom audits of exhaustive representation, nonsurjectivity, the omitted
real theorem and the realization-relation obstruction use only
\texttt{propext}, \texttt{Classical.choice} and \texttt{Quot.sound}.
No admitted proofs, additional axioms, native evaluation shortcuts or
auxiliary computation are used. The complete graph model includes all
finite sizes and arbitrary grounding; labeling artifacts or drawing
geometry cannot add real-valued game information.
\end{document}
